\subsection{A matching query law for the carrier family}
\label{sec:pn-family-query-law}

The preceding lower construction also admits a matching upper bound
on input probes. This result concerns that explicit family on its
entire input sign cube. It does not identify the optimal exponent
for arbitrary pre-norm networks.

\begin{theorem}[The carrier family's query exponent]
\label{thm:pn-family-query-law}
Fix a legal dyadic $s>\tanh(1)/\sqrt2$. Consider the networks
constructed in Theorem~\ref{thm:pn-normalized-gain}, allowing every
initial feature, carrier, and inert coordinate to be an arbitrary
sign. Let $N$ be the largest power of two at most $n/2$, put
$\eta=N/n$, $r_D=1+\sum_{j<D}a_j$, and define
\[
 \nu_\eta=\frac{s}{\tanh(1)\sqrt\eta}-1.
\]
Their optimal worst-input expected query cost is
\[
 \Theta_s\!\left(\min\{r_D^{2\nu_\eta},n\}\right).
\]
An exact sampler attaining the upper bound has finite expected
bit work. A direct certified implementation uses
\[
 O_s\!\left(
 \min\{(D+1)r_D^{2\nu_\eta},n+D+1\}\,B^4\right),
\]
where $B\ge16$ includes all parameter and radius-prefix encoding
lengths and $\log(n+D+2)$. Theorem~\ref{thm:pn-family-bit-law}
removes the depth factor by preprocessing the two-mean function.
\end{theorem}

\paragraph{Proof idea.}
On the full cube, the network depends on two unknown averages: the
feature mean and the carrier mean. Its normalized scalar state
has derivative growth $r_D^{\nu_\eta}$. The discrepancy from that
limiting gain has summable contribution along the step schedule.
Derivatives through order six then fit the two-variable corrected
Bernstein factory with a base sample size of order
$r_D^{2\nu_\eta}$. One observed feature sign determines which of
two bounded head functions the factory must simulate.

\begin{proof}
Write $c=\tanh1$, and let $z$ and $w$ denote the initial feature
and carrier means. Inert coordinates always have squared value one.
Every feature coordinate is $x_i-z+F_j(z,w)$. The squared RMS
denominator is
\[
 v_j=4-\eta z^2+\eta F_j^2-2\eta c\tau_j w
                         +\eta c^2\tau_j^2,
 \qquad \tau_j=r_j-1.
\]
Set $u_j=F_j/r_j$ and $q_j=a_j/r_{j+1}$. Then
\[
 u_{j+1}=(1-q_j)u_j+q_j T_j(z,w,u_j),\qquad
 T_j(z,w,u)=\tanh\!\left(
 \frac{su}{\sqrt{b_j(z,w)+\eta u^2}}\right),
\]
where
\[
 b_j(z,w)=
 \frac{4-\eta z^2-2\eta c\tau_j w+\eta c^2\tau_j^2}{r_j^2}.
\]
For all $(z,w)\in[-1,1]^2$ and $j$,
\[
 r_j^2b_j\ge4-2\eta+\eta(c\tau_j-1)^2
              \ge3+\eta(c\tau_j-1)^2.
\]
The carrier-floor identity in \eqref{eq:pn-carrier-floor}
therefore gives $b_j\ge1/12$. Moreover
\[
 |b_j-\eta c^2|\le8/r_j.
\]
For example, expansion in $1/r_j$ gives first-order coefficient
$-2\eta c(c+w)$ of magnitude at most two, and a second-order
coefficient of magnitude at most six. The state satisfies
$|u_j|\le1$, since its update is a convex combination of numbers
in $[-1,1]$.

Put $\nu=\nu_\eta$, $C_u=192s$, and $M=e^{C_u}$. Direct
differentiation and the mean-value theorem give
\[
 0\le\partial_uT_j
 =\operatorname{sech}^2\!\left(
       \frac{su}{\sqrt{b_j+\eta u^2}}\right)
       \frac{s b_j}{(b_j+\eta u^2)^{3/2}}
 \le1+\nu+C_u/r_j.
\]
Here the derivative of $x^{-1/2}$ on $x\ge1/12$ is at most
$12^{3/2}/2$, and $4\cdot12^{3/2}<192$.
Consequently a derivative of the scalar state propagated from
level $k$ to level $j$ has multiplier at most
\begin{equation}
 \prod_{i=k}^{j-1}(1+\nu q_i+C_u q_i/r_i)
 \le M(r_j/r_k)^\nu.
 \label{eq:pn-family-derivative-product}
\end{equation}
This uses $\sum q_i\le\log(r_j/r_k)$ and
$\sum q_i/r_i=1/r_k-1/r_j\le1$.

We give explicit uniform local derivative bounds. Set
\[
 h=\frac1{10000(1+s)},\qquad C_0=1440(3/h)^6.
\]
For real $z,w,u\in[-1,1]$, extend $T_j$ to the complex polydisk
of coordinate radius $h$. The squared denominator changes by
at most $4h$, while its real value is at least $1/12$.
The principal inverse square root is therefore holomorphic there.
Its argument inside tanh changes by at most
$6s h\,12^{3/2}<1/8$. Tanh has no pole there and has modulus
at most two. Cauchy's estimates show that the entrywise $\ell_1$
norm of every local derivative tensor of order at most six is
at most $C_0$; the factor $3^6$ covers the sum over coordinates.
The same bound holds for the head function, after enlarging $C_0$
by an absolute factor if necessary; the argument below permits
the choice $2C_0$ instead.

Let $E_k(j)$ be the maximum entrywise $\ell_1$ norm of the
$k$th derivative tensor of $u_j$ on $[-1,1]^2$. For $k=1$,
the direct parameter derivatives have total magnitude at most
$C_0$. Equations \eqref{eq:pn-family-derivative-product} and
\[
 \sum_iq_i r_{i+1}^{-\nu}\le1/\nu
\]
give $E_1(j)\le C_1r_j^\nu$, where
$C_1=M(1+C_0/\nu)$.
For $k=2,\ldots,6$, define recursively
\[
 H_k=\max\{2+C_1,C_2,\ldots,C_{k-1}\},\qquad
 C_k=\frac{M C_0 k^k H_k^k}{(k-1)\nu}.
\]
Then
\begin{equation}
 E_k(j)\le C_k r_j^{k\nu}.
 \label{eq:pn-family-derivatives}
\end{equation}
To verify this induction, use the multivariate chain rule in
partition form. The single-block term is
$\partial_uT_j\,D^ku_j$. Every other partition has at least
two blocks, all of size below $k$. There are at most $k^k$
partitions. The first derivative of the map
$(z,w)\mapsto(z,w,u_j)$ has total norm at most
$(2+C_1)r_j^\nu$; each higher derivative is bounded by the
induction hypothesis. Thus the remainder has norm at most
$C_0 k^kH_k^k r_j^{k\nu}$.
After propagation by \eqref{eq:pn-family-derivative-product},
its sum is bounded using
\[
 \sum_i q_i r_i^p\le(r_j^p-1)/p,
 \qquad p>0.
\]
For this inequality, integrate $t^{p-1}$ on
$[r_i,r_{i+1}]$ and use
$t^{p-1}\ge r_i^p/r_{i+1}$. This proves
\eqref{eq:pn-family-derivatives} with the displayed constants.

Read $x_1=e\in\{-1,1\}$ once. The required output-sign mean is
the known function
\[
 f_e(z,w)=\tanh\!\left(
 \frac{u_D(z,w)+(e-z)/r_D}
 {4\sqrt{b_D(z,w)+\eta u_D(z,w)^2}}\right).
\]
The full-domain head estimate in
Theorem~\ref{thm:pn-normalized-gain} gives
$|f_e|\le\tanh(1/\sqrt3)<4/5$, even at parameter pairs
inconsistent with the observed sign. The same local complex
argument just used applies to the head: its numerator has
absolute value at most three, its squared denominator is at
least $1/12$, and changing each variable by $h$ changes its
tanh argument by less than $1/8$.
The partition chain rule and \eqref{eq:pn-family-derivatives}
therefore bound its derivatives of orders two through six by
$A_{k,s,\eta}r_D^{k\nu}$, with explicit constants obtained
from the displayed $C_k$ and $2C_0 k^k
\max\{2+C_1,C_2,\ldots,C_k\}^k$.

Convert the means to Bernoulli biases by
$p_1=(1+z)/2$, $p_2=(1+w)/2$; this multiplies the order-$k$
derivative bounds by $2^k$. The two-variable corrected Bernstein
lemma, Lemma~\ref{lem:rank-two-bernstein}, requires known rational
derivative budgets. Enclose
$k\nu\log_2 r_D$ by certified logarithm and exponential
arithmetic, and round each displayed derivative bound upward
to a dyadic rational within a fixed factor. This costs
$O_s(B^4)$ work and gives budgets of $O_s(B)$ encoding length.
The lemma applies with base size $m_0=O_{s,\eta}(r_D^{2\nu})$.
Each source bit draws a uniformly random feature or carrier
coordinate, using at most one input probe. The two streams use
independent coordinate draws. They are independent conditional
on every fixed input, including the already observed $x_1$.
The returned sign has exactly the desired mean. Alternatively,
read all $2N$ relevant coordinates and evaluate the two known
means directly. This proves the upper query bound.

For clarity, the coefficient computations can be implemented
directly without an uncharged function oracle. Evaluate the
scalar recurrence and its bivariate derivatives through order
four by the chain rule, using certified dyadic interval arithmetic.
There are $O(D+1)$ fixed-size derivative operations. Local
derivative coefficients are bounded by the constants above.
Here is an explicit form of the error estimate used for guards.
Let $e_k(j)$ bound the error in the order-$k$ derivative tensor,
for $0\le k\le4$, and suppose every local update has absolute
rounding error at most $\epsilon$. Clip only the approximate
state value to $[-1,1]$, which cannot increase its error.
On the known derivative ranges, the product rule gives
\[
 e_k(j+1)\le(1+\nu q_j+C_uq_j/r_j)e_k(j)
 +C_s q_j\sum_{l<k}r_j^{(k-l)\nu}e_l(j)+\epsilon.
\]
For $k=0$ the sum is empty. The constants are obtained by
differentiating the finite partition formulas; they use local
derivatives through order five, already bounded above. Induction
on $k$, followed by \eqref{eq:pn-family-derivative-product} and
the integral estimate with $p=k\nu$, gives
\[
 e_k(j)\le A_{k,s}(D+1)\epsilon r_j^{(k+1)\nu}.
\]
The rounding-error sum contributes at most
$M(D+1)\epsilon r_j^\nu$; each lower-order forcing term has
the displayed exponent $(k+1)\nu$, which proves this claim.
The same estimate applies to the head derivatives. Choosing
$\epsilon\le2^{-t}/[C_s(D+1)r_D^{5\nu}]$ both gives the
requested final accuracy and closes the assumption that the
computed derivatives lie within one of their known bounds.
Evaluating the fixed number of local operations to an additional
$O_s(\log(r_D+2))$ guard bits makes their aggregate error at most
$\epsilon$. The resulting guard length is
$O_s(\log(D+2)+\log(r_D+2))$, already included in $O_s(B)$.
The known constant $c$ is enclosed to the same accuracy. This
gives $O_s((D+1)(B+t+\operatorname{bits}(p_1,p_2))^4)$
bit work per value or derivative evaluation.

A Bernstein correction at scale $m$ requires one scalar
hypergeometric sum of $O(m)$ such evaluations. Its normalized
coefficient needs $O_s(B+\log(m+2))$ initial precision; further
precision has a geometric tail under the fresh uniform comparison.
The correction masses are $O_s(m_0^2/m^2)$ on dyadic levels,
so summing the resulting expected costs gives
$O_s((D+1)m_0B^4)$. Reading all inputs instead costs
$O_s((n+D+1)B^4)$. These two implementations give the displayed
finite bit bound, including all scalar decisions and refinement.

For the lower bound, repeat
Theorem~\ref{thm:pn-normalized-gain} on the subdomain with all
carrier signs equal to $-1$, now retaining $A_0=\eta c^2$
and $A_1=4+2\eta c(1-c)$ in its normalized-gain comparison.
This gives exponent $2\nu_\eta$ with the same fixed readout
and a full-cube floor at least $1/12$.
Finally $\eta\in[1/4,1/2]$ and
$\nu_\eta\ge\sqrt2s/c-1>0$. All constants in both proofs
are bounded uniformly on this compact parameter interval.
For example, use $\eta c^2\ge c^2/4$,
$\nu_\eta\ge\sqrt2s/c-1$, and
$\nu_\eta\le2s/c-1$ in every prefactor. Thus the constants
may be chosen to depend only on $s$, as claimed.
\end{proof}
